\section{Demo Flows And Helper Scripts}
\label{sec:DemoFlows}

Durga is easier to understand when treated as something to run and probe, not only
as a generator to inspect statically. For that reason, generated helper scripts are
part of the intended product surface.

\subsection{Purpose of the helper layer}

The generated scripts serve three distinct roles:
\begin{itemize}
\item they provide runnable examples of the generated Kafka topology,
\item they reduce the amount of manual Kafka producer and consumer work needed during exploration,
\item they make BPMN semantics visible through explicit actions such as completion, failure, escalation and observation.
\end{itemize}

This matters because many workflow designs seem plausible in BPMN notation, but only become
clear once concrete runtime events are observed. Durga therefore generates a probe layer
alongside the executable runtime.

\subsection{Common generated scripts}

The exact set depends on the BPMN model, but the common script categories are:
\begin{itemize}
\item \texttt{demo-scenario.sh} for whole-process happy, stuck or failed runs,
\item \texttt{send-task-input.sh} for manual injection into one task input topic,
\item \texttt{complete-task.sh}, \texttt{fail-task.sh} and \texttt{escalate-task.sh} for human or exceptional task outcomes,
\item \texttt{complete-call-activity.sh} for returning from a generated call activity,
\item \texttt{send-message-event.sh} and \texttt{send-signal-event.sh} for external event stimulation,
\item \texttt{watch-process-events.sh} and \texttt{watch-task-output.sh} for runtime observation.
\end{itemize}

\subsection{Recommended learning loop}

The most effective way to understand a generated process is usually:
\begin{enumerate}
\item generate a project from one small BPMN sample,
\item create the Kafka topics and start the generated runtime,
\item run \texttt{watch-process-events.sh},
\item execute \texttt{demo-scenario.sh happy},
\item repeat with a failure, escalation or timeout-oriented script,
\item inspect the monitoring endpoints or dashboard after each run.
\end{enumerate}

This loop exposes the difference between nominal flow and exceptional flow with very little setup.
It also makes one design principle visible: the execution topics carry work, but the canonical
\texttt{process-events} stream carries the lifecycle truth.

\subsection{A practical division of use}

In practice the helper surface is most useful in three situations:
\begin{description}
\item[Learning] A user wants to understand how one BPMN construct maps to Kafka consumers, producers and events.
\item[Testing] A developer wants to force a boundary, event-subprocess or call-activity path without creating a custom producer.
\item[Debugging] A generated process compiles, but the user needs to see where runtime routing or cancellation deviates from the BPMN intent.
\end{description}

\subsection{Why this is documented in the manual}

This documentation is deliberate. The helper scripts are not incidental wrappers around the runtime.
They are part of how Durga explains itself. A BPMN-to-Kafka tool with no easy way to trigger and
observe the generated process would be much harder to evaluate, especially when exploring tradeoffs
around timers, boundaries, event subprocesses and subprocess scopes.
